\documentclass[letterpaper,11pt]{article}

\usepackage{latexsym}
\usepackage[empty]{fullpage}
\usepackage{titlesec}
\usepackage{marvosym}
\usepackage[usenames,dvipsnames]{color}
\usepackage{verbatim}
\usepackage{enumitem}
\usepackage[hidelinks]{hyperref}
\usepackage{fancyhdr}
\usepackage[english]{babel}
\usepackage{tabularx}
\input{glyphtounicode}

\pagestyle{fancy}
\fancyhf{}
\fancyfoot{}
\renewcommand{\headrulewidth}{0pt}
\renewcommand{\footrulewidth}{0pt}

\addtolength{\oddsidemargin}{-0.5in}
\addtolength{\evensidemargin}{-0.5in}
\addtolength{\textwidth}{1in}
\addtolength{\topmargin}{-.5in}
\addtolength{\textheight}{1.0in}

\urlstyle{same}
\raggedbottom
\raggedright
\setlength{\tabcolsep}{0in}

\titleformat{\section}{
  \vspace{-4pt}\scshape\raggedright\large
}{}{0em}{}[\color{black}\titlerule \vspace{-5pt}]

\pdfgentounicode=1

%-------------------------
% Custom commands

\newcommand{\resumeItem}[1]{
  \item\small{
    {#1 \vspace{-2pt}}
  }
}

\newcommand{\resumeSubheading}[4]{
  \vspace{0pt}\item
    \begin{tabular*}{0.97\textwidth}[t]{l@{\extracolsep{\fill}}r}
      \textbf{#1} & #2 \\
      \textit{\small #3} & \\
    \end{tabular*}
    \textit{\small #4}
    \vspace{-7pt}
}

\newcommand{\resumeProjectHeading}[2]{
  \item
    \begin{tabular*}{0.97\textwidth}{l@{\extracolsep{\fill}}r}
      \small #1 & #2 \\
    \end{tabular*}\vspace{-7pt}
}

\renewcommand\labelitemii{$\vcenter{\hbox{\tiny$\bullet$}}$}

\newcommand{\resumeSubHeadingListStart}{\begin{itemize}[leftmargin=0.10in, label={}]}
\newcommand{\resumeSubHeadingListEnd}{\end{itemize}}
\newcommand{\resumeItemListStart}{\begin{itemize}[leftmargin=0.14in]}
\newcommand{\resumeItemListEnd}{\end{itemize}\vspace{-8pt}}

%-------------------------------------------
% RESUME

\begin{document}

%----------HEADING----------

\begin{center}
    \textbf{\Huge \scshape Sharzil Hasnine} \\ \vspace{5pt}
    +61 451 645 516 $|$
    \href{mailto:sharzilhasnine975@gmail.com}{sharzilhasnine975@gmail.com} $|$
    \href{https://github.com/Shas-9}{github.com/Shas-9} \\
    Portfolio: \href{https://shas-9.github.io/}{shas-9.github.io}
\end{center}

%-----------EDUCATION-----------

\section{Education}
\resumeSubHeadingListStart

\resumeSubheading
  {Bachelor of Computer Science (Advanced)}
  {\small Feb 2023 -- Nov 2026}
  {Adelaide University}
  {}

\resumeSubHeadingListEnd

%-----------TECHNICAL SKILLS-----------




%-----------TECHNICAL EXPERIENCE-----------

\section{Experience}
  \resumeSubHeadingListStart

   \resumeSubheading
         {\textbf{{Industry Project - Block Model Compression Algorithm - Maptek}}}{\small Aug 2025 - Oct 2025}
         {\emph{\textbf{Tech:}} C++,  Bash, Jira}{}
         \resumeItemListStart
           \resumeItem{Architected a stream-oriented lossless compression system for 3D data, optimising for speed and compression ratio.}
           \resumeItem{Implemented an octree-based spatial compression algorithm, achieving ~98.8\% reduction in data size while improving processing speed by 2.6x over the uncompressed baseline, forming the foundation for subsequent team optimisations reaching 99.17\% compression and 1520x processing speed.}
           \resumeItem{Designed a testing framework that verifies correctness of compressed outputs against uncompressed input data, ensuring algorithmic correctness independent of compression strategy and enabling automated, extensible testing.}
           \resumeItem{Built a cross-platform CI pipeline with automated testing and reproducible Linux and Windows builds, integrating cross-compilation and artifact generation into the development workflow, improving team velocity by 1.5x.}
         \resumeItemListEnd

    \resumeSubheading{\textbf{Technical Volunteer - Autonomous navigation - Adelaide Rover Team}}{\small Oct 2024 - Feb 2025}
    {\emph{\textbf{Tech:}} Python, C++, Ros2, Nav2, Gazebo, Groot}{}
    \resumeItemListStart
        \resumeItem{Designed a behavior tree for the rover’s autonomous navigation pipeline using Nav2, enabling modular decision-making for tasks such as path planning, recovery behaviours, and goal execution.}
        \resumeItem{Simulated navigation behaviour using Gazebo, supporting the Adelaide Rover Team's preparation for the Australian Rover Challenge's autonomous navigation task.}
    \resumeItemListEnd
    

  \resumeSubHeadingListEnd


%-----------TECHNICAL PROJECTS-----------

\section{Projects}

\resumeSubHeadingListStart

\resumeSubheading
  {CatQuest -- 2D Maze Exploration Game}
  {\small Aug 2023 -- Oct 2023}
  {Tech: C++, SFML}
  {\textbf{Link:} \href{https://github.com/Shas-9/Labyrinth-game}{github.com/Shas-9/Labyrinth-game}}

\resumeItemListStart

\resumeItem{
Developed gameplay systems for a real-time 2D maze exploration game in C++ using SFML, working with the main game loop, player state and interactive entities.
}

\resumeItem{
Implemented enemy behaviour, player mechanics, collision detection and health systems while integrating them with an existing procedurally generated maze.
}

\resumeItem{
Designed modular game-state and UI architecture to support menu navigation and transitions between different gameplay states without disrupting the real-time game loop.
}

\resumeItem{
Debugged interactions between gameplay systems and refined implementation details to maintain predictable behaviour during real-time execution.
}

\resumeItemListEnd

\resumeSubheading
  {Monkey Run -- Unihack 2024}
  {\small Mar 2024}
  {Tech: Godot, GDScript, Python, MediaPipe}
  {\textbf{Link:} \href{https://github.com/rayokamoto/unihack}{github.com/rayokamoto/unihack}}

\resumeItemListStart

\resumeItem{
Developed an endless-runner game in Godot using GDScript, including gameplay logic and randomised obstacle generation.
}

\resumeItem{
Implemented a Python and MediaPipe input system that translated real-time webcam-tracked player movement into game controls.
}

\resumeItem{
Integrated the external input system with gameplay mechanics, debugging interactions between multiple components to deliver a responsive interactive experience.
}

\resumeItemListEnd

\resumeSubHeadingListEnd


\end{document}